\begin{remark*}[Exposition]
The interval operations on $\mathcal{I}(\mathbb{R})$ preserve some familiar
algebraic laws from real arithmetic, but not all of them. Addition and
multiplication remain commutative and associative, and the singleton intervals
$[0,0]$ and $[1,1]$ serve as neutral elements. However, nondegenerate intervals
generally do not have additive or multiplicative inverses, and the ordinary
distributive law fails in general.
\end{remark*}

\begin{proposition}[Commutativity and associativity of interval addition and multiplication]
\cite{UnknownIntroductionIntervalAnalysis2002}
\label{prop:interval-addition-multiplication-commutative-associative}
For $A,B,C\in\mathcal{I}(\mathbb{R})$,
\[
A+B=B+A,
\qquad
(A+B)+C=A+(B+C),
\]
and
\[
A\cdot B=B\cdot A,
\qquad
(A\cdot B)\cdot C=A\cdot(B\cdot C).
\]
\end{proposition}

\begin{proposition}[Neutral intervals]
\cite{UnknownIntroductionIntervalAnalysis2002}
\label{prop:neutral-intervals}
For every $A\in\mathcal{I}(\mathbb{R})$,
\[
A+[0,0]=[0,0]+A=A
\]
and
\[
A\cdot[1,1]=[1,1]\cdot A=A.
\]
\end{proposition}

\begin{remark*}[Nonexistence of inverses]
A nondegenerate interval has no inverse with respect to interval addition or
interval multiplication, because combining a nondegenerate interval with another
interval cannot produce a singleton neutral interval.
\end{remark*}

\begin{example}[Failure of distributivity]
\cite{UnknownIntroductionIntervalAnalysis2002}
\label{ex:failure-of-interval-distributivity}
The distributive law does not hold in general for interval arithmetic. For instance,
\[
[1,2]\cdot\bigl([1,2]-[1,2]\bigr)
=
[1,2]\cdot[-1,1]
=
[-2,2],
\]
whereas
\[
[1,2]\cdot[1,2]-[1,2]\cdot[1,2]
=
[1,4]-[1,4]
=
[-3,3].
\]
Thus
\[
[1,2]\cdot\bigl([1,2]-[1,2]\bigr)
\neq
[1,2]\cdot[1,2]-[1,2]\cdot[1,2].
\]
\end{example}

\begin{proposition}[Sub-distributivity of interval multiplication]
\cite{UnknownIntroductionIntervalAnalysis2002}
\label{prop:interval-subdistributivity}
For $A,B,C\in\mathcal{I}(\mathbb{R})$,
\[
A\cdot(B+C)\subseteq A\cdot B+A\cdot C.
\]
\end{proposition}

\begin{remark*}[Special distributive case]
The ordinary distributive rule is valid when the multiplying interval is a
singleton real scalar. In particular, for $t\in\mathbb{R}$,
\[
t(A+B)=tA+tB.
\]
\end{remark*}
